\uuid{UxfU}
\exo7id{4922}
\titre{exo7 4922}
\auteur{quercia}
\organisation{exo7}
\datecreate{2010-03-17}
\isIndication{false}
\isCorrection{true}
\chapitre{Conique}
\sousChapitre{Hyperbole}
\cpge{mpsi-mp2i.calcul-algebrique-trigonometrie.d, pcsi.calcul-algebrique-trigonometrie.d, ptsi.calcul-algebrique-trigonometrie.d, mpsi-mp2i.espaces-prehilbertiens.b, pcsi.espaces-prehilbertiens.b, ptsi.geometrie-espace.b, bcpst1.trigonometrie, pt.courbes-implicites-plan.b, pt.structure-prehilbertienne.b}
\module{Géométrie}
\niveau{L2}
\difficulte{}

\contenu{
\question{
Soient $O,A$ deux points distincts du plan.
Trouver les points $M$ tels que
$\overline{(\vec{OA},\vec{OM})} \equiv 2\overline{(\vec{AM},\vec{AO})}$.
}
\reponse{
Soit $\alpha \equiv \overline{(\vec{AM},\vec{AO})}$.\par
         $\frac{OM}{\sin\alpha} = \frac{AM}{\sin(2\alpha)}
           \Rightarrow  2\cos\alpha = \frac{AM}{OM}$.\par
         Al-Khâshi $ \Rightarrow  \frac{AM^2}{OM}(OM-OA) = (OM-OA)(OM+OA)$.\par
         $OM = OA  \Rightarrow  \alpha \equiv \frac\pi4$.\par
         $OM \ne OA  \Rightarrow  OM = 2d(M,\Delta)$ où $\Delta$ est la médiatrice de
         $[O,A]$, et $M$ est du côté de $O$.
}
}
